% Complete Section 4.4, original PDF 49-54 / printed 39-44.
% The user-approved non-Fraktur mappings preserve harmonic/cohomology distinctions.
% block: abstract-hodge-heading page:49 type:section
\subsection{{抽象 Hodge Laplacian}}\label{sec:abstract-hodge-laplacian}

% block: hodge-laplacian-domain-intro page:49 type:ordinary-paragraph
{设 $(W,d)$ 为闭 Hilbert 复形. 与之相伴的是由 $L=dd^*+d^*d$ 给出的次数为 $0$ 的分次算子 $W\to W$. 对与 Riemannian 流形相伴的 de Rham 复形, 此算子称为 Hodge Laplacian; 一般地, 我们将它称为与该复形相伴的 (抽象) Hodge Laplacian. 具体地,}
\begin{equation}\label{eq:abstract-hodge-domain}
D(L^k)=\{\,u\in V^k\cap V_k^*\mid du\in V_{k+1}^*,\ d^*u\in V^{k-1}\,\},
\end{equation}
{且}
\begin{equation}\label{eq:abstract-hodge-operator}
\begin{gathered}
L^k u=d_{k+1}^*d^k u+d^{k-1}d_k^*u=(d^k)^*d^k u+d^{k-1}(d^{k-1})^*u,\qquad u\in D(L^k).
\end{gathered}
\end{equation}

% block: hodge-kernel-identification page:49 type:ordinary-paragraph
{FEEC 的一个主要课题是对给定 $f\in W^k$, 数值求解 Hodge Laplace 方程 $L^k u=f$. 这一问题是否具有唯一解, 取决于调和形式空间 $\mathsf H^k$ 是否为零. 事实上, 下面证明 $L^k$ 的零空间恰好为 $\mathsf H^k$, 由此解释\emph{调和}这一术语. 由 \eqref{eq:harmonic-cocycle-orthogonality} 可知 $\mathsf H^k\subset D(L^k)$, 且对所有 $u\in\mathsf H^k$ 有 $L^k u=0$. 因此 $\mathsf H^k\subset\mathcal N(L^k)$. 另一方面, 由 $L$ 的定义有}
\begin{equation}\label{eq:hodge-green-identity}
\langle L^k u,v\rangle=\langle du,dv\rangle+\langle d^*u,d^*v\rangle,\qquad u\in D(L^k),\ v\in V^k\cap V_k^*.
\end{equation}
{取 $v=u$, 得}
\[
\|du\|^2+\|d^*u\|^2=\langle L^k u,u\rangle,\qquad u\in D(L^k),
\]
{因而,}
\[
\|du\|^2+\|d^*u\|^2=0,
\]
{只要 $u\in\mathcal N(L^k)$. 因此, 若 $u\in\mathcal N(L^k)$, 则 $du$ 与 $d^*u$ 均为零, 从而 $u\in\mathsf H^k$. 这验证了 $\mathcal N(L^k)=\mathsf H^k$, 即 Hodge Laplacian 的解只在相差一个调和形式的意义下确定.}

% block: hodge-solvability-orthogonality page:49 type:ordinary-paragraph
{调和形式还决定了解存在的一个必要条件. 将方程 $L^k u=f$ 与 $q\in\mathsf H^k$ 作内积, 并利用 \eqref{eq:hodge-green-identity}, 得 $\langle f,q\rangle=0$. 因此, 解存在的一个必要条件是正交条件 $f\perp\mathsf H^k$. 事实上, 下面的 Theorem~\ref{thm:hodge-well-posedness} 将证明这一条件不仅必要, 而且足以保证解的存在.}

% block: hodge-three-formulations-heading page:50 type:subsection
\subsubsection{{Hodge Laplace 问题的三种表述}}\label{subsec:hodge-three-formulations}

% block: hodge-strong-form-intro page:50 type:ordinary-paragraph
{为了构造对任意 $f\in W^k$ 都可解的问题, 利用 $W^k$ 到 $\mathsf H^k$ 上的正交投影 $P_{\mathsf H}$. 由此可提出 Hodge Laplace 方程的如下\emph{强表述}, 我们将证明它是适定的.}

% block: hodge-strong-formulation page:50 type:ordinary-paragraph
\medskip\noindent\textbf{{强表述}}\par
{给定 $f\in W^k$, 求 $u\in D(L^k)$, 使得}
\begin{equation}\label{eq:hodge-strong-formulation}
Lu=f-P_{\mathsf H}f,\qquad u\perp\mathsf H^k.
\end{equation}

% block: hodge-solution-decomposition page:50 type:ordinary-paragraph
{注意, 若 $u$ 为此问题的一个解, 则}
\begin{equation}\label{eq:hodge-data-decomposition}
f=dd^*u+P_{\mathsf H}f+d^*du.
\end{equation}
{显然 $dd^*u\in\mathcal B$, $P_{\mathsf H}f\in\mathsf H$, $d^*du\in\mathcal B^*$, 因而 \eqref{eq:hodge-data-decomposition} 是 $f$ 的 Hodge 分解, 即 $dd^*u=P_{\mathcal B}f$, $d^*du=P_{\mathcal B^*}f$.}

% block: hodge-weak-formulations-intro page:50 type:ordinary-paragraph
{还将考虑该问题的两种弱表述. 原始表述直接对应于强表述.}

% block: hodge-primal-formulation page:50 type:ordinary-paragraph
\medskip\noindent\textbf{{原始弱表述}}\par
{给定 $f\in W^k$, 求 $u\in V^k\cap V_k^*$, 使 $u\perp\mathsf H^k$ 且}
\begin{equation}\label{eq:hodge-primal-formulation}
\langle du,dv\rangle+\langle d^*u,d^*v\rangle=\langle f-P_{\mathsf H}f,v\rangle,\qquad v\in V^k\cap V_k^*.
\end{equation}

% block: hodge-mixed-variables page:50 type:ordinary-paragraph
{对混合弱表述, 引入两个新变量 $\sigma=d^*u\in V^{k-1}$ 和 $p=P_{\mathsf H}f\in\mathsf H^k$.}

% block: hodge-mixed-formulation page:50 type:ordinary-paragraph
\medskip\noindent\textbf{{混合弱表述}}\par
{给定 $f\in W^k$, 求 $\sigma\in V^{k-1}$, $u\in V^k$, $p\in\mathsf H^k$, 使得}
\begin{equation}\label{eq:hodge-mixed-formulation}
\begin{aligned}
\langle\sigma,\tau\rangle-\langle u,d\tau\rangle&=0,&\tau&\in V^{k-1},\\
\langle d\sigma,v\rangle+\langle du,dv\rangle+\langle p,v\rangle&=\langle f,v\rangle,&v&\in V^k,\\
\langle u,q\rangle&=0,&q&\in\mathsf H^k.
\end{aligned}
\end{equation}

% block: hodge-mixed-asymmetry page:50 type:ordinary-paragraph
{一个重要观察是, 与另两种表述不同, 混合弱表述并不对称地处理算子 $d$ 与 $d^*$. 这种表述中只出现算子 $d$, 涉及 $d^*$ 的项以弱形式处理.}

% block: hodge-formulation-equivalence-intro page:50 type:ordinary-paragraph
{尽管形式不同, 这三种表述完全等价.}

% block: hodge-formulation-equivalence page:50 type:theorem
\Needspace{8\baselineskip}
\begin{Theorem}{{表述的等价性}}\label{thm:hodge-formulation-equivalence}
{设 $f\in W^k$. 元素 $u\in W^k$ 是抽象 Hodge Laplace 问题强表述的解, 当且仅当它是原始弱表述的解. 此外, 此时若令 $\sigma=d^*u$, $p=P_{\mathsf H}f$, 则三元组 $(\sigma,u,p)$ 是混合弱表述的解. 最后, 若某个 $(\sigma,u,p)$ 是混合弱表述的解, 则 $u$ 是该问题另两种表述的解, 且 $\sigma=d^*u$, $p=P_{\mathsf H}f$.}
\end{Theorem}

% block: hodge-equivalence-proof-regularity page:51 type:proof
\begin{Proof}
{先证明强表述与原始弱表述的等价性. 若 $u$ 满足强表述, 则当然有 $u\in V\cap V^*$ 且满足原始弱表述. 另一方面, 若 $u$ 满足原始弱表述, 并且\emph{若 $u$ 属于 $D(L)$}, 则它满足强表述. 因此, 需证明任意满足原始弱表述, 从而先验地只属于 $V\cap V^*$ 的 $u$, 实际上属于 \eqref{eq:abstract-hodge-domain} 定义的 $D(L)$, 即满足 $du\in V^*$ 与 $d^*u\in V$. $du\in V^*$ 意味着存在常数 $c$, 使得}
\begin{equation}\label{eq:hodge-du-adjoint-bound}
\langle du,dv\rangle\leq c\|v\|,\qquad v\in V.
\end{equation}

% block: hodge-equivalence-proof-hodge-split page:51 type:proof
{利用 Hodge 分解, 将 $v=v_1+v_2$ 分解, 其中 $v_1\in\mathcal Z$, $v_2\in V\cap\mathcal B^*$. 显然,}
\[
\langle du,dv_1\rangle=0,
\]
{而}
\[
\langle du,dv_2\rangle=\langle du,dv_2\rangle+\langle d^*u,d^*v_2\rangle=\langle f-P_{\mathsf H}f,v_2\rangle\leq c\|v_2\|\leq c\|v\|,
\]
{其中 $c=\|f-P_{\mathsf H}f\|$. 将这些结果相加得到 \eqref{eq:hodge-du-adjoint-bound}, 从而表明 $du\in V^*$. 类似的论证表明 $d^*u\in V$, 因而证明了 $u\in D(L)$, 也就证明原始弱表述的解是强解.}

% block: hodge-equivalence-proof-strong-mixed page:51 type:proof
{其次, 设 $u$ 为强表述的解, 并令 $\sigma=d^*u$, $p=P_{\mathsf H}f$. 由于 $u\in D(L)$, 有 $\sigma\in V$. 混合表述的第一个方程由 $\sigma$ 的定义推出. 由 Hodge Laplace 方程有}
\[
d\sigma+d^*du+p=f,
\]
{与任意 $v\in V$ 作内积, 得到混合方法的第二个方程. 第三个方程由假设 $u\perp\mathsf H^k$ 得到. 这表明 $(\sigma,u,p)$ 是弱表述的解.}

% block: hodge-equivalence-proof-mixed-primal page:51 type:proof
{最后, 若 $(\sigma,u,p)\in V\times V\times\mathsf H$ 是混合弱表述的解, 则第一个方程给出 $u\in V^*$ 和 $d^*u=\sigma$, 最后一个方程给出 $u\perp\mathsf H$. 在第二个方程中取 $v\in\mathsf H$, 可知 $p=P_{\mathsf H}f$. 最后, 在第二个方程中取 $v\in V\cap V^*$, 可得 $u$ 是原始弱表述的解.}
\end{Proof}

% block: hodge-well-posedness-heading page:51 type:subsection
\subsubsection{{适定性}}\label{subsec:hodge-well-posedness}

% block: hodge-well-posedness-intro page:51 type:ordinary-paragraph
{最后证明, 在计入调和形式之后, 抽象 Hodge Laplace 方程是适定的.}

% block: hodge-well-posedness page:51 type:theorem
\Needspace{8\baselineskip}
\begin{Theorem}{{适定性}}\label{thm:hodge-well-posedness}
{对每个 $f\in W^k$, 存在唯一解 $u\in D(L^k)$, 满足 $u\perp\mathsf H^k$ 且 $Lu=f-P_{\mathsf H}f$. 此外,}
\[
\|u\|+\|du\|+\|d^*u\|+\|dd^*u\|+\|d^*du\|+\|p\|\leq c\|f\|.
\]
{该估计中的常数 $c$ 仅依赖于 Poincaré 不等式 (Theorem~\ref{thm:hilbert-complex-poincare}) 中的常数 $c_P$.}
\end{Theorem}

% block: hodge-well-posedness-proof-reduction page:52 type:proof
\begin{Proof}
{当解存在时, $\|dd^*u\|$, $\|d^*du\|$ 和 $\|p\|$ 的界直接来自 Hodge 分解 $f=dd^*u+d^*du+p$, 因此只须证明对任意 $f\in W$, 存在唯一解且满足}
\[
\|u\|_V+\|d^*u\|\leq c\|f\|.
\]

% block: mixed-hodge-bilinear-form page:52 type:proof
{证明基于混合弱表述. 简记 $X=V^{k-1}\times V^k\times\mathsf H^k$, 将定义混合表述的三个方程的左端合并, 定义有界双线性型 $B:X\times X\to\mathbb R$:}
\begin{equation}\label{eq:hodge-mixed-bilinear-form}
\begin{aligned}
B(\sigma,u,p;\tau,v,q)={}&\langle\sigma,\tau\rangle-\langle u,d\tau\rangle-\langle d\sigma,v\rangle\\
&-\langle du,dv\rangle-\langle p,v\rangle-\langle u,q\rangle.
\end{aligned}
\end{equation}

% block: mixed-hodge-load-functional page:52 type:proof
{注意 $B$ 是对称的. 还通过下式定义有界线性泛函 $F:X\to\mathbb R$:}
\begin{equation}\label{eq:hodge-mixed-load-functional}
F(\tau,v,q)=\langle f,v\rangle.
\end{equation}
{则混合弱表述具有标准形式: 求 $x\in X$, 使得}
\begin{equation}\label{eq:hodge-mixed-abstract-problem}
B(x,y)=F(y),\qquad y\in X.
\end{equation}

% block: hodge-necas-inf-sup page:52 type:proof
{Nečas 对 Lax–Milgram 定理的推广 [\BookNumber{61}, Theorem 3.1] 给出了这一问题何时适定. (对任意 $F\in X'$ 的) 适定性等价于 \emph{inf-sup 条件}}
\begin{equation}\label{eq:hodge-inf-sup-condition}
\gamma:=\inf_{0\ne w\in X}\sup_{0\ne y\in X}\frac{B(w,y)}{\|w\|_X\|y\|_X}>0,
\end{equation}
{此时解 $x$ 满足 $\|x\|_X\leq\gamma^{-1}\|F\|_{X'}$. (一般还需要交换 $w$ 与 $y$ 角色后的第二个 inf-sup 条件, 但本情形中 $B$ 是对称的, 因而不需要.) 下面的定理证明这一 inf-sup 条件.}
\end{Proof}

% block: hodge-inf-sup-theorem page:52 type:theorem
\Needspace{6\baselineskip}
\begin{Theorem}\label{thm:hodge-mixed-inf-sup}
{\eqref{eq:hodge-mixed-bilinear-form} 定义的双线性型 $B$ 满足 inf-sup 条件 \eqref{eq:hodge-inf-sup-condition}, 其中 $\gamma$ 的下界仅依赖于 Poincaré 常数 $c_P$.}
\end{Theorem}

% block: hodge-inf-sup-proof-task page:52 type:proof
\begin{Proof}
{对任意给定的 $(\sigma,u,p)\in V^{k-1}\times V^k\times\mathsf H$, 须构造不恒为零的 $(\tau,v,q)\in V^{k-1}\times V^k\times\mathsf H$, 使得}
\[
B(\sigma,u,p;\tau,v,q)\geq\gamma(\|\sigma\|_V+\|u\|_V+\|p\|)(\|\tau\|_V+\|v\|_V+\|q\|),
\]
{其中 $\gamma$ 仅依赖于 $c_P$. 为此, 使用 Hodge 分解 $u=u_{\mathcal B}+u_{\mathsf H}+u_{\mathcal B^*}$, 其中 $u_{\mathcal B}=P_{\mathcal B}u$, $u_{\mathsf H}=P_{\mathsf H}u$, $u_{\mathcal B^*}=P_{\mathcal B^*}u$. 由于 $u_{\mathcal B}\in\mathcal B^k$, 可写成 $u_{\mathcal B}=d\rho$, 其中 $\rho\in\mathcal Z^{k-1\perp}$. 由 Poincaré 不等式,}
\begin{equation}\label{eq:hodge-inf-sup-lifting-bound}
\|\rho\|_V\leq c_P\|u_{\mathcal B}\|.
\end{equation}

% block: hodge-inf-sup-test-choice page:52 type:proof
{然后定义}
\begin{equation}\label{eq:hodge-inf-sup-test-choice}
\begin{gathered}
\tau=\sigma-\frac{1}{c_P^2}\rho\in V^{k-1},\qquad v=-u-d\sigma-p\in V^k,\\
q=-p+u_{\mathsf H}\in\mathsf H^k.
\end{gathered}
\end{equation}

% block: hodge-inf-sup-test-bound page:53 type:proof
{由 \eqref{eq:hodge-inf-sup-lifting-bound} 及 Hodge 分解的正交性, 有}
\begin{equation}\label{eq:hodge-inf-sup-test-bound}
\|\tau\|_V+\|v\|_V+\|q\|\leq C(\|\sigma\|_V+\|u\|_V+\|p\|).
\end{equation}
{将 \eqref{eq:hodge-inf-sup-test-choice} 代入 \eqref{eq:hodge-mixed-bilinear-form}, 得}
\[
\begin{aligned}
B(\sigma,u,p;\tau,v,q)={}&\|\sigma\|^2+\|d\sigma\|^2+\|du\|^2+\|p\|^2+\|u_{\mathsf H}\|^2\\
&+\frac{1}{c_P^2}\|u_{\mathcal B}\|^2-\frac{1}{c_P^2}\langle\sigma,\rho\rangle.
\end{aligned}
\]

% block: hodge-inf-sup-cross-term page:53 type:proof
{利用 Cauchy–Schwarz 不等式, 代数-几何平均不等式以及 \eqref{eq:hodge-inf-sup-lifting-bound} 估计最后一项:}
\[
\begin{aligned}
\langle\sigma,\rho\rangle&\leq\|\sigma\|\|\rho\|\leq\frac{c_P^2}{2}\|\sigma\|^2+\frac{1}{2c_P^2}\|\rho\|^2\\
&\leq\frac{c_P^2}{2}\|\sigma\|^2+\frac12\|u_{\mathcal B}\|^2.
\end{aligned}
\]
{因此}
\[
\begin{aligned}
B(\sigma,u,p;\tau,v,q)\geq{}&\frac12\|\sigma\|^2+\|d\sigma\|^2+\|du\|^2+\|p\|^2\\
&+\|u_{\mathsf H}\|^2+\frac{1}{2c_P^2}\|u_{\mathcal B}\|^2.
\end{aligned}
\]

% block: hodge-inf-sup-dual-boundary-bound page:53 type:proof
{现对 $u_{\mathcal B^*}$ 应用 Poincaré 不等式, 其中 $du_{\mathcal B^*}=du$:}
\[
\|u_{\mathcal B^*}\|_V\leq c_P\|du\|.
\]
{得到}
\[
\begin{aligned}
&B(\sigma,u,p;\tau,v,q)\\
&\geq\frac12\|\sigma\|^2+\|d\sigma\|^2+\frac12\|du\|^2+\|p\|^2+\|u_{\mathsf H}\|^2\\
&\qquad+\frac{1}{2c_P^2}\|u_{\mathcal B}\|^2+\frac{1}{2c_P^2}\|u_{\mathcal B^*}\|^2\\
&\geq\eta(\|\sigma\|^2+\|d\sigma\|^2+\|u_{\mathcal B}\|^2+\|u_{\mathsf H}\|^2+\|u_{\mathcal B^*}\|^2+\|du\|^2+\|p\|^2)\\
&=\eta(\|\sigma\|_V^2+\|u\|_V^2+\|p\|^2),
\end{aligned}
\]
{其中 $\eta=\min(1/2,1/2c_P^2)$ 仅依赖于 $c_P$. 由此界及 \eqref{eq:hodge-inf-sup-test-bound}, 容易得到 inf-sup 条件.}
\end{Proof}

% block: hodge-general-load-remark page:53 type:remark
\Needspace{6\baselineskip}
\begin{Remark}\label{rem:hodge-general-mixed-load}
{Theorem~\ref{thm:hodge-mixed-inf-sup} 中的 inf-sup 条件意味着混合问题 \eqref{eq:hodge-mixed-abstract-problem} 对一般 $F\in X'$ 适定, 而不仅对具有 \eqref{eq:hodge-mixed-load-functional} 形式的 $F$ 适定.}
\end{Remark}

% block: hodge-solution-operator page:53 type:ordinary-paragraph
{已确立 Hodge Laplacian 的适定性, 将解算子记为 $K$. 即对 $f\in W$, 通过方程 $LKf=f-P_{\mathsf H}f$, $Kf\perp\mathsf H$ 定义 $Kf\in D(L)$. 用 $K$ 表示, 混合表述的解由 $\sigma=d^*Kf$, $u=Kf$, $p=P_{\mathsf H}f$ 给出. 因而 Hodge 投影为 $P_{\mathcal B}=dd^*K$, $P_{\mathcal B^*}=d^*dK$, 且}
\[
u=dd^*Ku+P_{\mathsf H}u+d^*dKu,\qquad u\in W.
\]

% block: hodge-mixed-choice-remark page:53 type:remark
\Needspace{8\baselineskip}
\begin{Remark}\label{rem:hodge-mixed-discretization-motivation}
{读者可能会问, 为什么选择用混合表述而非原始弱表述来确立抽象 Hodge Laplacian 的适定性. 毕竟, 原始表述的双线性型由 Poincaré 不等式可知是强制的, 因此可以应用 Riesz 表示定理, 无须验证 inf-sup 条件. 原因在于我们着眼于 Hodge Laplacian 的离散化, 这将是下一章的主题. 离散化将基于混合表述, 正如将要看到的, 混合表述远比原始表述更适于离散化. 混合双线性型的 inf-sup 条件 (而不仅是对具有 \eqref{eq:hodge-mixed-load-functional} 形式的 $F$ 的适定性) 届时将成为离散化方法分析的关键.}
\end{Remark}

% block: boundary-problems-heading page:54 type:subsection
\subsubsection{{$\mathcal B$ 与 $\mathcal B^*$ 问题}}\label{subsec:boundary-coboundary-problems}

% block: boundary-problem-hodge-reduction page:54 type:ordinary-paragraph
{本小节考察数据 $g$ 属于 $\mathcal B$ 或 $\mathcal B^*$ 时的 Hodge Laplace 方程 $Lv=g$. 前一情形中, 对某个 $f\in V$ 有 $g=df$ (若 $g\in V^k$, 则 $f\in V^{k-1}$). 由 Hodge Laplace 方程}
\begin{equation}\label{eq:hodge-boundary-problem-potential}
Lu=f-P_{\mathsf H}f
\end{equation}
{在数据为 $f$ 时的解 $u$, 令 $v=du$ 即得到 $Lv=g$ 的解. 这是合理的, 因为形式上有}
\[
Lv=(dd^*+d^*d)du=dd^*du=d(dd^*+d^*d)u=d(f-P_{\mathsf H}f)=g.
\]

% block: boundary-problem-domain-justification page:54 type:ordinary-paragraph
{只要 $v=du\in D(L)$, 这一形式计算就有依据. 为说明这一点, 注意由于 $u\in D(L)$, 有 $v=du\in V^*$ 且 $v\in V$. 此外 $dv=0$, 并且由 $f\in V$ 及 $\mathcal Z\subset V$, 有 $d^*v=d^*du=P_{\mathcal B^*}f=f-P_{\mathcal Z}f\in V$. 这就表明 $v\in D(L)$, 如前所述.}

% block: boundary-problem-intro page:54 type:ordinary-paragraph
{因此, 当 Hodge Laplacian 的数据 $g$ 属于 $\mathcal B$ 时, 相应解 $v$ 也属于 $\mathcal B$. 此时 Hodge Laplacian 中的项 $d^*dv$ 为零, 所以 $v$ 满足另一个问题.}

% block: boundary-problem-formulation page:54 type:ordinary-paragraph
\medskip\noindent\textbf{{$\mathcal B$ 问题}}\par
{给定 $g\in\mathcal B^k$, 求 $v\in\mathcal B^k$, 使 $dd^*v=g$.}

% block: boundary-problem-equivalence page:54 type:ordinary-paragraph
{按定义, 此问题的解还必须属于 $D(dd^*)=\{\,v\in V_k^*\mid d^*v\in V^{k-1}\,\}$. 刚才已看到, 对 $g\in\mathcal B^k$, 若 $v$ 是数据为 $g$ 的 Hodge Laplace 问题的解, 则 $v$ 是 $\mathcal B$ 问题的解. 反之, 若 $v$ 是 $\mathcal B$ 问题的解, 则有 $dv=0$, 从而 $Lv=dd^*v=g$, 即 $v$ 是 Hodge Laplace 问题的解. 由 Hodge Laplace 问题的适定性, 得到 $\mathcal B$ 问题的适定性 (但仅对 $g\in\mathcal B^k$; 否则无解).}

% block: boundary-problem-well-posedness page:54 type:theorem
\Needspace{6\baselineskip}
\begin{Theorem}\label{thm:hodge-boundary-problem}
{若 $g\in\mathcal B^k$, 则数据为 $g$ 的 $\mathcal B$ 问题具有唯一解 $v$, 且此解与相同数据的 Hodge Laplace 问题的唯一解相同.}
\end{Theorem}

% block: dual-boundary-problem-intro page:54 type:ordinary-paragraph
{当然, 对 $d^*$ 也有类似结论: 若 $f\in V^*$, 且 $u$ 为相应 Hodge Laplace 问题的解, 则 $v=d^*u\in D(L)$, 并且它是数据为 $g=d^*f$ 的 Hodge Laplace 问题的解. 此外, 若 $g\in\mathcal B^*$, 则它也是下述问题的唯一解.}

% block: dual-boundary-problem-formulation page:54 type:ordinary-paragraph
\medskip\noindent\textbf{{$\mathcal B^*$ 问题}}\par
{给定 $g\in\mathcal B_k^*$, 求 $v\in\mathcal B_k^*$, 使 $d^*dv=g$.}
